\documentclass[11pt, letterpaper]{article}

\topmargin 0in \textheight 9in \textwidth 6.7in \evensidemargin -.21in \oddsidemargin -.21in \columnsep 0.3in

\usepackage{amsmath}
\usepackage{amsfonts}
\usepackage{amssymb}
\usepackage{graphicx}
\usepackage{array}
\usepackage{booktabs}
\usepackage{enumitem}
\usepackage{hyperref}
\usepackage{xcolor}
\usepackage{natbib}
\usepackage[T1]{fontenc}
\usepackage{titlesec}
\titleformat{\section}{\normalfont\normalsize\bfseries}{}{0pt}{}
\titlespacing*{\section}{0pt}{10pt}{4pt}
\usepackage{fancyhdr}
\pagestyle{fancy}
\fancyhf{}
\renewcommand{\headrulewidth}{0pt}
\fancyfoot[L]{\footnotesize Proposal length: 2 pages max, not including references.}
\fancyfoot[R]{\footnotesize\thepage}

\hypersetup{
    colorlinks=true,
    linkcolor=blue!70!black,
    urlcolor=blue!70!black,
    citecolor=blue!70!black
}

\setlength{\parindent}{0pt}
\setlength{\parskip}{6pt}


\begin{document}
{\Large\textbf{ECE750T: Research Project Proposal}}\\[5pt]
{\large\textbf{Your Project Title}}\\[4pt]
Your Name \quad | \quad Your Student ID\\
University of Waterloo \quad | \quad Fall 2026
\vspace{0.5em}\hrule\vspace{0.5em}

% Replace all bracketed text. Keep the eight headings.
% At most 2 pages excluding references. This template is for the proposal only.
% Released September 25, 2026. Consult the handout for the submission deadline.

\section*{Q1: Research Question}
[State your question in plain language and explain its relevance to Responsible AI (possibly including topics we haven't yet covered in class).]

\section*{Q2: Prior Work and Limitations}
[Describe 3--5 relevant works and their limitations. Cite the papers you discuss.
Add your own entries to references.bib and cite them with
\texttt{\textbackslash citep\{key\}} or \texttt{\textbackslash citet\{key\}}.
Example citation (replace with your own sources): \citet{chouldechova2017fair}.]

\section*{Q3: Novelty of Approach}
[Explain what is new or different about your approach, application, or comparison.]

\section*{Q4: Impact and Motivation}
[Explain why the question matters and who would benefit from the results.]

\section*{Q5: Risks and Assumptions}
[Identify key assumptions and possible failure modes.]

\section*{Q6: Resources}
[Describe the data, compute, or other materials you need and how you will obtain
them. Distinguish resources already available from access still to be arranged.]

\newpage
\section*{Q7: Weekly Plan}
% The draft planning window is October 9--December 4, 2026.
% Update these dates if the confirmed course deadlines change.
\begin{center}
\renewcommand{\arraystretch}{1.2}
\begin{tabular}{lp{4.3in}}
\toprule
\textbf{Week} & \textbf{Plan} \\
\midrule
Oct 12--18 (reading week) & [Milestone] \\
Oct 19--25 & [Milestone] \\
Oct 26--Nov 1 & [Milestone] \\
Nov 2--8 & [Milestone] \\
Nov 9--15 & [Milestone] \\
Nov 16--22 & [Milestone] \\
Nov 23--29 & [Milestone] \\
Nov 30--Dec 4 & [Milestone] \\
\bottomrule
\end{tabular}
\end{center}

\section*{Q8: Milestones and Success Criteria}
[Define evaluation criteria, baselines or other comparisons, and what you expect
to learn. Include quantitative and/or qualitative criteria as appropriate.]

\section*{Tools and Assistance}
[If you plan to use AI tools in your research, briefly describe how. As relevant, mention any other relevant resources (archives, software used within your lab) that will play a significant role in your project's success.]

% References are excluded from the two-page limit. Replace the example citation
% and bibliography entry with sources relevant to your proposal.
% latexmk runs BibTeX automatically.
\bibliographystyle{plainnat}
\bibliography{references}
\end{document}
